Stochastic Gradient Descent (SGD) is the core optimization tool driving modern machine learning. Recent years have seen substantial progress in understanding its dynamics, particularly in two-layer networks~\citep{saad1995line,mei2018mean,chizat2018global,rotskoff2022trainability,sirignano2020mean,arnaboldi2023high}. While global convergence is qualitatively well-understood when the network is wide enough, quantitative results are scarcer. A particularly fruitful body of recent theoretical work addressing this gap has focused on deriving precise convergence rates for particular model classes on synthetic data, such as high-dimensional  Gaussian single and multi-index models~\citep{arous2021online,abbe2022merged,abbe2023sgd}. These advances have sparked a wave of follow-up studies~\citep{damian2022neural,damian2023smoothing,bietti2023learning,ba2023learning,moniri2023theory,mousavihosseini2023gradient,cui2024asymptotics,zweig2023symmetric,berthier2023learning,arnaboldi2024escaping,arnaboldi2024online}, deepening our understanding of what problems are hard to learn for neural networks trained under SGD.


While multi-index models have served as a cornerstone in theoretical analyses of learning, they remain far from the architectures driving recent breakthroughs in machine learning. Modern advances in learning from sequential data --- particularly in natural language processing --- are increasingly dominated by attention-based models such as the Transformer architecture~\citep{vaswani2017attention}. These models introduce a paradigm shift through self-attention mechanisms, which dynamically reweight the influence of each input token based on its relevance to others. Through successive layers of attention, Transformers capture intricate dependencies across sequences, enabling state-of-the-art performance in tasks ranging from machine translation to large-scale language modeling~\citep{brown2020language,kenton2019bert}.

The main goal of this work is to extend our theoretical understanding of learning with SGD on multi-index models to attention-based architectures and sequential data. Inspired by~\cite{cui2024phase,troiani2025fundamental}, our focus will be on the following class of single-layer, tied attention models:
\begin{equation} \label{2506:eq:model-definition}
    \vec{f}_{\vec{w}}(X) = \Pi\left[
        \SoftMax{\left(\left(X+\frac{P}{\sqrt{d}}\right)\vec{w}\vec{w}^\top\left(X+\frac{P}{\sqrt{d}}\right)^\top\right)}
    \right],
\end{equation}
where $\vec{w}\in \mathbb{R}^{d}$ are the trainable weights and $\Pi$ is the \emph{reduction map}, which allows passing from an $L\times L$ matrix to a $k_\mathrm{out}$-dimensional vector. As detailed in~\Cref{2506:app:reduction}, this model is a reduction of the standard attention mechanism, where
\begin{enumerate}
    \item key and query matrix are tied, and $d_\mathrm{head} = 1$: $Q=K=X\cdot \vec{w}\in \mathbb{R}^{L\times 1}$ ; 
    \item since we are considering a single layer attention, and we do not need to learn a new representation of the sequence, the value matrix is the identity: $V=I_{L}$.
\end{enumerate}

In this work, we will be interested in the optimization properties of the model in~\cref{2506:eq:model-definition} when trained under (spherical) one-pass stochastic gradient descent (SGD) from a random initial condition $\vec{w}^0 \sim {\rm Unif}(\mathbb{S}^{d-1}(\sqrt{d}))$:
\begin{equation}\label{2506:eq:normalized-SGD-w}
\begin{split}
    \vec{w}^{\tau+1} &= \frac{\vec{w}^\tau - \gamma \nabla_{\vec{w}} \ell(X^\tau, \vec{y}^\tau; f_{\vec{w}})}{\left\lVert \vec{w}^\tau - \gamma \nabla_{\vec{w}} \ell(X^\tau, \vec{y}^\tau; f_{\vec{w}})\right\rVert}
    \norm{\vec{w}^\tau}.
\end{split}\end{equation}
with the squared loss:
\begin{equation}
    \ell(X, \vec{y}; f_{\vec{w}}) =
        \left\lVert\vec{y} - f_{\vec{w}}(X)\right\rVert^2_{\rm F} =
        \sum_{i=1}^{k_\mathrm{out}} \left(y_i - f_{\vec{w}}(X)_i\right)^2.
\end{equation}
Note that for one-pass (a.k.a. \emph{online} or \emph{streaming}) SGD each sample is only seen once, meaning that the sample complexity of the algorithm coincides with the convergence rate. The spherical constraint is considered to simplify the mathematical analysis, a common assumption in the analysis of SGD for single-index models~\citep{arous2021online,damian2022neural}.

To derive a sharp characterization of the sample complexity and convergence rate of SGD for the single-layer attention mechanism in~\cref{2506:eq:model-definition}, we assume that the sequence data $(X, \vec{y})$ is generated from the following Gaussian \emph{sequence single-index} (SSI) model:
\begin{assumption}[Data distribution] 
\label{2506:def:gsim}
We assume training data $(X, \vec{y})\in\mathbb{R}^{L\times d}\times \mathbb{R}^{k_\mathrm{out}}$ is independently drawn from a Gaussian \emph{Sequence Single Index (SSI)} model:
    \begin{equation}
    \label{2506:eq:gssim}
        \vec{f}^\mathrm{SSI}_{\vec w_{\star}}(X) = \vec{g}{\left(X\cdot \vec{w}_{\star}\right)}
    \end{equation}
    where $X\in\mathbb{R}^{L\times d}$ is a Gaussian matrix with entries $\mathcal{N}(0,\sfrac{1}{d})$, \(\vec{g}\colon \mathbb{R}^L\to\mathbb{R}^{k_\mathrm{out}}\) is a vector-valued link function that depends only on the scalar product of the input with a fixed vector \(\vec{w}_\star\in\mathbb{S}^{d-1}(\sqrt{d})\) and \(k_\mathrm{out}\) is an integer that does not depend on \(d\) or \(L\).  
\end{assumption}  
Recent work by~\citep{cui2024phase,troiani2025fundamental,cui2025high} has shown that the single-layer tied attention model in~\cref{2506:eq:model-definition} is a particular instance of a class of \emph{sequence multi-index (SMI) models}, creating a bridge between single-index analysis and attention-based learning. This mapping implies that model~\cref{2506:eq:model-definition} can learn at best a predictor in this class, justifying the choice for training data. The SSI model class was introduced by~\cite{cui2024phase} in the context of studying phase transitions for the attention mechanism. These results position the SMI model as a signature synthetic model to study the interplay between attention mechanisms, data structure, and the dynamics of learning algorithms. Despite this progress, the learning dynamics of SMI models under SGD remain unexplored. Prior studies~\citep{cui2025high,cui2024phase} analyzed the empirical risk minimizer through the heuristic replica method, while~\citep{troiani2025fundamental} rigorously studied the Bayes-optimal estimator for SMI models. However, a theoretical understanding of the population landscape and SGD dynamics in this model is missing. Our work addresses this gap, providing the first rigorous characterization of SGD dynamics in SMI models by leveraging techniques developed for single and multi-index models. 


\paragraph{Main results ---} Our main methodological contribution is the generalization of analytical tools to study multi-index models to variants that process sequences of tokens rather than simple vector inputs. Our contributions can be summarized as follows:
\begin{itemize}
    \item We introduce the notion of \emph{sequence information exponent} (SIE), as a generalization of the information exponent for single-index models~\cite{arous2021online}; the SIE has a direct correspondence with the sample complexity of SGD. We also discuss the implications of the positional encoding on the sample complexity, proving it could help SGD to learn faster.
    \item We analyze the speed-up introduced by the attention mechanism when learning sequential data, compared to models  not adapted to treat sequence structures, e.g. fully connected networks. For many problems, we show that the gain is proportional to the sequence length \(L\) and in some cases even larger. 
    \item We investigate the interplay between the positional  and semantic structure of the data following the setting from~\cite{cui2024phase}, showing that SGD dynamics is not always able to disentangle the two and that a rich phase diagram arises describing the structure of the corresponding population loss and the performance of SGD. 
\end{itemize}
All our formal claims are supported by rigorous proofs, as well as numerical experiments; the code developed is available at this \href{https://github.com/IdePHICS/Sequence-Single-Index}{GitHub repository}.

\paragraph{Further Related works ---} There has been much activity discussing {\bf SGD with synthetic data on multi-index models} over the last decades or so, see e.g.~\citep{arous2021online,veiga2022phase,arnaboldi2023high,collinswoodfin2023hitting,marion2023leveraging,berthier2023learning,montanari2025dynamical} and references therein. {\bf Information} and  {\bf generative} exponent have been the topic of many works over the last few years~\citep{arous2021online,abbe2022merged,damian2024computational,troiani2024fundamental,defilippis2025optimal}. Here we discuss and adapt these notions for sequence models~\cite{troiani2025fundamental}.

On the theory of transformers, \cite{sanford2023representational} has shown that transformers can approximate certain sparse averaging (qSA) tasks independently of the sequence length, while fully connected networks require a width scaling with the length. \cite{mousavihosseini2025when} derived sample complexity upper bounds for q-Sparse Token Regression (qSTR) tasks, a random version of qSA. \cite{wang2024transformers} showed that the separation between transformers and fully-connected networks also hold under GD for the qSTR task. While the qSA and qSTR tasks are not directly comparable to the SSIM, both model a sparse dependency of the task in the sequence. \cite{marion2025attention} introduced a model that can be seen as a sequence two-index model. While we assume the input data to be i.i.d. Gaussian they assume a spiked covariance structure in the data which makes their results not directly comparable to ours. 

Convergence guarantees for a single layer transformer were studied by~\cite{wu2023convergence}. \cite{song2024unraveling} studied the convergence of GD for simple architectures and highlighted the existence of suboptimal local solutions. 
\cite{li2025optimization} point out that rapid convergence does not guarantee meaningful learning.  \cite{li2023theoretical} give sample complexity bounds for a shallow vision transformer. \cite{zhang2025understanding} studied benign overfitting in transformers trained by SGD. \cite{yuksel2025sample} derived sample complexity upper bounds for next token prediction tasks based on the complexity of the hypothesis class, and showed that the out-of-sample error depends on the sequence mixing properties. Compared to these work we move beyond convergence results to study how the data structure, e.g. sequence structure and positional encodings, affect sample complexity and behaviors of the population loss and recovery dynamics in the high-dimensional limit.
